\bA_T=\bG_T\herm\bG_T-\lambda_{\rm null}\bH_u\bH_u\herm, \qquad
\lambda_{\rm null}=0.5, \label{eq:wT}
\end{equation}
followed by unit-norm normalization. This trades target illumination against
communication leakage, consistent with leakage-aware ISAC beam design
\cite{wang2025hybrid,lyu2024crb}.

With $g_{kj}=|\bh_k\herm\bw_j|^2$, the SINRs are
\begin{align}
\gamma_f &= \frac{a_fg_{ff}P}{a_ng_{ff}P+a_Tg_{fT}P+\sigma^2}, \label{eq:gf}\\
\gamma_n &= \frac{a_ng_{nn}P}{a_Tg_{nT}P+\sigma^2}, \label{eq:gn}\\
\gamma_s &= \frac{a_T|\bm{u}_T\herm\bG_T\bw_T|^2P}{\sigma^2}, \label{eq:gs}
\end{align}
where $\sigma^2$ is the noise power and $\bm{u}_T$ is the sensing receive
combiner. The achievable rates are $R_k=\log_2(1+\gamma_k)$,
$k\in\{n,f,s\}$.

The weighted system rate is
\begin{equation}
\Rdl=\omega_c(R_n+R_f)+\omega_sR_s,\qquad \omega_c+\omega_s=1.
\label{eq:rdl}
\end{equation}
The communication rates are summed in the \emph{objective} because the near and
far users occupy the same NOMA resource block, so $R_n+R_f$ is the cluster
spectral efficiency. The weights $\omega_c$ and $\omega_s$ are design
priorities that trade communication throughput against sensing performance; in
our simulations they are fixed to $0.7$ and $0.3$, respectively. Importantly,
using a sum in the objective does not replace individual QoS protection.

\subsection{Problem Formulation}\label{sec:prob}
We jointly optimize the power split $\bm{a}$ and RIS phases
$\bm{\phi}=[\phi_1,\ldots,\phi_N]$. The communication sum rate is retained as a
cluster-level QoS measure, while individual user constraints are added in the
same spirit as the per-device QoS constraint in (9b) of~\cite{zou2023ml}. The
resulting problem is
\begin{subequations}\label{eq:prob}
\begin{align}
\text{(P1)}\quad \max_{\bm{a},\bm{\phi}}\quad & \Rdl \label{eq:obj}\\
\text{s.t.}\quad
& R_n+R_f\geq R_{\mathrm{th},c}, \label{eq:qosc}\\
& R_k\geq R_{\mathrm{th},k},\quad k\in\{n,f\}, \label{eq:qosu}\\
& R_s\geq R_{\mathrm{th},s}, \label{eq:qoss}\\
& a_n+a_f+a_T=1,\quad a_n,a_f,a_T\geq0, \label{eq:simplex}\\
& |e^{j\phi_i}|=1,\quad i=1,\ldots,N. \label{eq:unit}
\end{align}
\end{subequations}
The aggregate constraint~\eqref{eq:qosc} is useful because the two NOMA users
share one time-frequency resource block, so $R_n+R_f$ directly measures the
communication service delivered by the cluster. It also matches the
cluster-level QoS term used by the implemented training loss. On its own,
however, a sum threshold is insufficient: a very large near-user rate could
compensate for an unacceptable far-user rate. Constraint~\eqref{eq:qosu}
therefore adds the individual protection requested in the formulation, exactly
analogous to writing a QoS condition for each device as in (9b) of
\cite{zou2023ml}. Thus, the sum threshold controls aggregate communication
quality while the individual thresholds prevent user starvation. The sensing
requirement is enforced separately by~\eqref{eq:qoss}. The ideal-RIS assumption
$\beta_i=1$ is fixed rather than optimized, and the phase parameterization
$e^{j\phi_i}$ enforces unit modulus automatically.

Problem~(P1) remains non-convex because the power variables and RIS phases are
coupled through the cascaded channels. We therefore enforce the QoS conditions
through penalty terms during learning, while the softmax output enforces the
power simplex exactly.

%==============================================================================
\section{Proposed MAML-Based Solution}
The implementation follows a bilevel power/phase decomposition. The
three-dimensional power vector is predicted by a compact neural network, while
the $N$ RIS phases are adapted for each channel realization. This follows the
basic motivation of~\cite{zou2023ml}: learn the lower-dimensional resource map
and refine the channel-specific RIS configuration with gradient-based
adaptation.

\subsection{Feature Map and Power-Allocation Network}\label{sec:feat}
For a given phase vector, the cascaded channels and the six effective gains
$\{g_{nn},g_{ff},g_{sT},g_{fn},g_{fT},g_{nT}\}$ are computed. They are converted
to dB and concatenated with the communication/sensing QoS levels and the SNR
budget, giving a nine-dimensional input feature vector $\bm z$. The power
network contains three hidden layers with 256 neurons per layer, LayerNorm and
LeakyReLU activations, followed by a three-neuron softmax head. Hence
\begin{equation}
\bm a=\operatorname{softmax}(f_{\bm\theta}(\bm z)),
\end{equation}
which enforces $a_n+a_f+a_T=1$ and $a_k\geq0$ by construction.

\subsection{Joint Power-Allocation and Phase Adaptation}
For a channel sample and the current phase vector $\bm\phi^{(j)}$, the network
first predicts
\begin{equation}
\bm a^{(j)}=\operatorname{softmax}
\left(f_{\bm\theta}(\bm z(\bm\phi^{(j)}))\right). \label{eq:powerpred}
\end{equation}
The implemented constrained loss is
\begin{align}
\ell^{(j)}={}&-\Rdl
+\lambda_c\big(R_{\mathrm{th},c}-R_n-R_f\big)_+
+\lambda_s\big(R_{\mathrm{th},s}-R_s\big)_+ \notag\\
&+\lambda_n\big(R_{\mathrm{th},n}-R_n\big)_+
+\lambda_f\big(R_{\mathrm{th},f}-R_f\big)_+, \label{eq:sampleloss}
\end{align}
where $(x)_+=\max(x,0)$. Equation~\eqref{eq:sampleloss} makes explicit why the
aggregate and individual communication thresholds coexist: $\lambda_c$ protects
cluster throughput, whereas $\lambda_n$ and $\lambda_f$ penalize starvation of
either user.

The RIS phase is represented by
$\theta_i=e^{j\phi_i}=\cos\phi_i+j\sin\phi_i$, so unit modulus is satisfied for
any real $\phi_i$. For each channel, the phase vector is randomly initialized
and adapted for $K$ steps using Adam on~\eqref{eq:sampleloss},
\begin{equation}
\bm\phi^{(j+1)}=\operatorname{Adam}_{\phi}
\left(\bm\phi^{(j)},\nabla_{\bm\phi^{(j)}}\ell^{(j)};\eta_{\phi}\right),
\label{eq:inner}
\end{equation}
with inner learning rate $\eta_{\phi}$. Since the effective gains change after
every phase update, the feature vector and power split are recomputed at each
inner step.

\subsection{First-Order Meta-Update}
After the $K$ phase-adaptation steps, the final rates and loss are recomputed at
$\bm\phi^{(K)}$. The implementation uses a first-order MAML approximation: the
adapted phase is treated as fixed during the outer network update rather than
backpropagating second-order derivatives through the entire Adam trajectory.
The outer update is therefore
\begin{equation}
\bm\theta\leftarrow\bm\theta-\alpha
\nabla_{\bm\theta}\frac{1}{B}\sum_{b=1}^{B}
\ell_b\big(\bm\theta,\bm\phi_b^{(K)}\big), \label{eq:outer}
\end{equation}
where $B$ is the mini-batch size. This retains the key meta-learning idea--the
power map is trained on channels after per-instance RIS adaptation--while
avoiding the memory and numerical burden of the full higher-order graph. In the
implemented fair, high-SINR configuration, the full second-order path was found
to produce unstable/non-finite meta-gradients, which motivated the first-order
variant.

\begin{figure}[t]
\centering
\begin{tikzpicture}[
  >=Stealth,
  node distance=5mm and 7mm,
  box/.style={draw,rounded corners,align=center,minimum width=20mm,minimum height=7mm,font=\scriptsize},
  arr/.style={->,thick},
  darr/.style={->,thick,dashed}]
\node[box,fill=yellow!12] (env) {Channel sample\\$\{\bG,\bm r_n,\bm r_f,\bG_T\}$};
\node[box,fill=blue!8,below left=of env] (phase) {RIS phase\\$\bm\phi^{(j)}$};
\node[box,fill=red!8,below right=of env] (net) {Power network\\$f_{\bm\theta}$};
\node[box,fill=green!10,below=12mm of env] (loss) {Rates + penalties\\$\ell^{(j)}$};
\draw[arr] (env) -- (phase);
\draw[arr] (env) -- (net);
\draw[arr] (phase) -- node[left,font=\scriptsize] {gains/features} (loss);
\draw[arr] (net) -- node[right,font=\scriptsize] {$\bm a^{(j)}$} (loss);
\draw[darr] (loss.west) to[bend left=18] node[below,font=\scriptsize] {Adam, $K$ steps} (phase.south);
\draw[darr] (loss.east) to[bend right=18] node[below,font=\scriptsize] {first-order outer update} (net.south);
\end{tikzpicture}
\caption{First-order MAML-based joint power/phase design. The inner loop adapts
the RIS phase and recomputes the network power split; the outer loop updates the
power network using the phase-adapted channel while ignoring higher-order
derivatives through the inner optimizer.}
\label{fig:maml}
\end{figure}

\subsection{Training and Prediction Procedure}
Algorithm~\ref{alg:maml} summarizes the implementation. At inference, the same
$K$-step phase adaptation is performed for an unseen channel and the final
softmax output provides the power allocation. Thus prediction is not a pure
one-shot neural-network call; it is a short learned-plus-optimization procedure
that trades a small fixed online cost for substantially better phase adaptation.

\begin{algorithm}[t]
\SetAlgoLined
\DontPrintSemicolon
\KwIn{channel data; inner steps $K$; inner Adam rate $\eta_{\phi}$; outer rate
$\alpha$; penalties $\lambda_c,\lambda_s,\lambda_n,\lambda_f$.}
\KwOut{trained power-map parameters $\bm\theta^\star$.}
Initialize $\bm\theta$\;
\While{not converged}{
 Sample a mini-batch of $B$ channels\;
 \ForEach{channel $b$}{
   Initialize $\bm\phi_b^{(0)}\sim\mathcal U([0,2\pi)^N)$\;
   \For{$j=0$ \KwTo $K-1$}{
     Compute phase-dependent channels/gains and $\bm z_b^{(j)}$\;
     Predict $\bm a_b^{(j)}$ using~\eqref{eq:powerpred}\;
     Evaluate $R_n,R_f,R_s$ and~\eqref{eq:sampleloss}\;
     Update $\bm\phi_b^{(j+1)}$ with inner Adam\;
   }
   Recompute $\bm a_b^{(K)}$ and final loss at the adapted phase\;
 }
 Update $\bm\theta$ with Adam using the first-order gradient in~\eqref{eq:outer}\;
}
Select $\bm\theta^\star$ by validation score\;
\caption{First-order MAML-based power and RIS-phase optimization}
\label{alg:maml}
\end{algorithm}

\subsection{Convergence Analysis}
A global convergence guarantee is not claimed. The objective is non-convex, the
unit-modulus RIS parameterization couples all rates through the cascaded channel,
and the hinge penalties in~\eqref{eq:sampleloss} are non-smooth exactly at the
QoS boundaries. Moreover, as discussed for the related RIS-NOMA MAML method in
\cite{zou2023ml}, standard smooth MAML convergence theorems do not directly
cover an architecture in which the inner and outer loops optimize different
variables. Our first-order approximation further discards the higher-order
dependence of the adapted phase on $\bm\theta$. Consequently, convergence is
assessed empirically: the inner Adam loop is run for a fixed number of steps and
the outer training curve is checked for a stable loss/rate plateau across the
reported hyperparameter settings.

\subsection{Computational Complexity}
Let $P_{\bm\theta}$ denote the power-network parameter count. With nine inputs,
three 256-unit hidden layers, LayerNorm, and a three-output head, the implemented
network contains approximately $1.36\times10^5$ trainable parameters. Computing
the user cascades costs $\mathcal O(MN)$, while forming the sensing matrix and
its $M\times M$ leakage-aware eigenvector adds approximately
$\mathcal O(MN+M^3)$. Therefore, one inner prediction/adaptation step costs
\begin{equation}
\mathcal O\!\left(P_{\bm\theta}+MN+M^3\right),
\end{equation}
and online prediction for one channel realization costs
